\chapter{النهايات والاتصال}\label{ch:b1:continuity}

استُعملت مبرهنة القيم الوسطى ومبرهنة القيم الحدّية
في مستوى الثانوي على الثقة البصرية. ومع المتتاليات
(\cref{ch:b1:seq}) وطوبولوجيا $\R$ (\cref{ch:b1:topology})
في اليد، يبرهن هذا الفصل عليهما --- ويكمّل النظرية
بمبرهنة التقابل الرتيب (التي تشرّع $\arcsin$ و
$\operatorname{arcosh}$ وأخواتهما من \cref{ch:b1:functions})
وبمبرهنة هاينه في \omterm{def:b1:continuity:uniform}{الاتصال المنتظم}.

وفي كل ما يأتي، $I$ \omterm{prop:b1:reals:intervals}{فترة} و $f \colon I \to \R$؛ و«$x_0 \in
\overline I$» يسمح بالنهايات عند الأطراف.

\section{نهايات الدوال}

\begin{definition}[النهاية عند نقطة]\label{def:b1:continuity:limit}
ليكن $x_0 \in \overline{I}$ و $\ell \in \R$. عندئذ يكون $f(x) \to \ell$ عندما
$x \to x_0$ إذا كان\index{نهاية!نهاية دالة}
\[
\forall \varepsilon > 0,\ \exists \delta > 0,\ \forall x \in I,
\qquad \abs{x - x_0} \leq \delta \implies \abs{f(x) - \ell} \leq
\varepsilon .
\]
وتُعرَّف النهايات عند $\pm\infty$ والنهايات
غير المنتهية بالنمط نفسه (إذ يصير $\abs{x - x_0} \leq \delta$ هو $x \geq M$؛ ويصير $\abs{f(x) -
\ell} \leq \varepsilon$ هو $f(x) \geq M'$). وتقصر النهايات من جهة واحدة
المتغيرَ $x$ على $x > x_0$ (ويُكتب $x \to x_0^+$) أو على $x < x_0$. والنهاية
وحيدة إذا وُجدت (بالبرهان نفسه في
\cref{prop:b1:seq:first}).
\end{definition}

\begin{example}[نهاية بالحصر]\label{ex:b1:continuity:floorlimit}
احسب $\lim_{x \to 0} x\,\bigl\lfloor \frac1x \bigr\rfloor$.
يعطي تأطير الجزء الصحيح $\frac1x - 1 < \lfloor \frac1x \rfloor \leq
\frac1x$، بعد الضرب في $x$ (وانتبه للإشارة!):
\[
1 - x < x\Bigl\lfloor \frac1x \Bigr\rfloor \leq 1 \quad (x > 0),
\qquad
1 \leq x\Bigl\lfloor \frac1x \Bigr\rfloor < 1 - x \quad (x < 0),
\]
وينغلق الحصران من الجهتين على $1$: فالنهاية هي $1$. ولاحظ
ما حدث: للمقدار $\lfloor \frac1x\rfloor$ وحده قفزات جامحة قرب
$0$، لكن العامل $x$ يروّض كل قفزة (فالمقدار $x$ مضروبًا في قفزة واحدية
صغير)، ولا ينجو إلا التأطير. والفكرة النافذة:
نهايات جداء عامل صغير في عامل ذي تذبذب
محدود مسائلُ حصر لا مسائل مبرهنة العمليات ---
إذ تحتاج مبرهنة العمليات إلى تقارب العاملين \emph{كليهما}.
\end{example}

\begin{theorem}[التمييز التتالي]\label{thm:b1:continuity:seqchar}
يكون $f(x) \to \ell$ عندما $x \to x_0$ إذا وفقط إذا: كان
$f(u_n) \to \ell$ من أجل \emph{كل} متتالية $(u_n)$ من نقاط $I$ مع $u_n \to x_0$.
\end{theorem}

\begin{proof}
($\Rightarrow$) ليكن $u_n \to x_0$ و $\varepsilon > 0$. خذ
$\delta$ من التعريف، ثم $N$ يحقق $\abs{u_n - x_0} \leq
\delta$ من أجل $n \geq N$: فبعد $N$، $\abs{f(u_n) - \ell} \leq
\varepsilon$.

($\Leftarrow$) بالنقيض. إذا كان $f \not\to \ell$: فإن
$\varepsilon_0 > 0$ ما يهزم كل $\delta$؛ ويعطي اختيار $\delta =
\frac{1}{n+1}$ عددًا $u_n \in I$ يحقق $\abs{u_n - x_0} \leq
\frac{1}{n+1}$ و $\abs{f(u_n) - \ell} > \varepsilon_0$. عندئذ $u_n
\to x_0$ لكن $f(u_n) \not\to \ell$.
\end{proof}

\begin{corollary}[العمليات والتركيب والترتيب]\label{cor:b1:continuity:operations}
تسلك مجاميع النهايات وجداءاتها وقسمتها (بنهاية غير معدومة في المقام)
سلوكَها من أجل المتتاليات؛ وإذا كان $f \to \ell$ عند $x_0$ و $g \to m$ عند $\ell$،
وكانت $g$ معرَّفة حول $\ell$ مع $g(\ell) = m$ أو كان $f \neq \ell$
بجوار $x_0$، فإن $g \circ f \to m$ عند $x_0$؛ وتحفظ النهايات
المتراجحات الواسعة، وتصحّ مبرهنة الحصر.
\end{corollary}

\begin{proof}
تنتقل كل \omterm{def:b1:logic:statement}{عبارة} عبر
\cref{thm:b1:continuity:seqchar} إلى نظيرتها التتالية
(\cref{thm:b1:seq:operations,thm:b1:seq:order}). وأمّا التركيب،
فالتسلسل المباشر يستحق أن يُكتب مرة: فمن أجل $\varepsilon >
0$ معطى، تعطي نهاية $g$ عند $\ell$ عددًا $\eta > 0$ يحقق
\[
\abs{y - \ell} \leq \eta \implies \abs{g(y) - m} \leq
\varepsilon \quad (y \text{ في مجموعة تعريف } g),
\]
حيث تُغطّى الحالة $y = \ell$ لأن $g(\ell) = m$؛ ثم
تعطي نهاية $f$ عند $x_0$ عددًا $\delta > 0$ يحقق $\abs{x -
x_0} \leq \delta \implies \abs{f(x) - \ell} \leq \eta$؛ وبالتسلسل،
يعطي $\abs{x - x_0} \leq \delta$ أن $\abs{g(f(x)) - m} \leq
\varepsilon$. وفي الفرض البديل (أي $f \neq \ell$ بجوار
$x_0$)، لا تُغذّى القيمة $y = \ell$ إلى $g$ أبدًا وقيمتها
هناك لا شأن لها.
\end{proof}

\begin{example}[لماذا يوجد شرط التركيب]\label{ex:b1:continuity:compositiontrap}
لتكن $g(y) = 0$ من أجل $y \neq 0$ و $g(0) = 1$، ولتكن $f$
مساوية $0$ تمامًا. عندئذ $f(x) \to 0$ عندما $x \to 0$، و $g(y) \to 0$
عندما $y \to 0$؛ ومع ذلك $g(f(x)) = g(0) = 1$ من أجل كل $x$، ومنه $g
\circ f \to 1 \neq 0$. فالدالة الداخلية تقف \emph{بالضبط على}
القيمة الممنوعة $\ell = 0$ إلى الأبد، وتتجاهل نهاية $g$ عند $0$
ما تفعله $g$ عند $0$. وشرط
\cref{cor:b1:continuity:operations} --- إمّا $g(\ell) = m$
(أي $g$ \omterm{def:b1:continuity:continuous}{متصلة} عند $\ell$)، وإمّا $f \neq \ell$ بجوار $x_0$ ---
هو بالضبط ما يستبعد هذا. والفكرة النافذة: في
الممارسة يركّب المرء دوالّ \emph{\omterm{def:b1:continuity:continuous}{متصلة}} ويكون
الشرط مجانيًا؛ وهو لا يعضّ إلا حين تُؤخذ النهايات على
جوارات مثقوبة، ولذلك يتضمن تعريف
$\lim_{x \to x_0}$ المستعمل في هذا الكتاب النقطةَ إذا كانت
في \omterm{def:b1:logic:sets}{مجموعة} التعريف.
\end{example}

\section{الاتصال}

\begin{definition}\label{def:b1:continuity:continuous}
تكون $f$ \emph{متصلة عند $x_0 \in I$}\index{اتصال} إذا كان $f(x)
\to f(x_0)$ عندما $x \to x_0$؛ و\emph{متصلة على $I$} إذا كانت
متصلة عند كل نقطة. وحسب
\cref{thm:b1:continuity:seqchar}: تكون $f$ متصلة عند $x_0$ إذا وفقط إذا كان
$f(u_n) \to f(x_0)$ من أجل كل متتالية $u_n \to x_0$ في $I$.
\end{definition}

\begin{example}[تصنيف الانقطاعات]\label{ex:b1:continuity:taxonomy}
ثلاث كيفيات للفشل عند نقطة، متزايدة الشدة.
\emph{قابل للرفع}: للمقدار $f(x) = \frac{\sin x}{x}$ على $\R^*$ نهايةٌ
$1$ عند $0$؛ ويصلحه تعريف $f(0) = 1$ --- فالانقطاع
كان ثقبًا لا سِمة. \emph{بقفزة}: للمقدار $\lfloor x \rfloor$ عند
عدد صحيح نهايتان متمايزتان من الجهتين ($n - 1$ و $n$)؛ ولا
اختيار للقيمة يوفّق بينهما، لكن نصفَي النهاية موجودان.
\emph{جوهري}: للمقدار $\sin\frac1x$ عند $0$ لا نهاية من أيّ جهة
أصلًا (\cref{exo:b1:continuity:1}) --- أي تذبذب دون
استقرار. وأمّا الدوال الرتيبة فلا تنتج إلا النوع الأوسط
(لأن نهاياتها من جهة واحدة موجودة دائمًا، لكونها حدودًا عليا ودنيا)،
ولذلك تكون مجموعات انقطاعها قابلة للعدّ على الأكثر ---
عدد ناطق واحد لكل قفزة. وأمّا المشتقات، فحسب مبرهنة دربو
(\cref{exo:b1:derivative:10})، لا تنتج إلا النوع الأخير:
فالدالة ذات الانقطاع بقفزة ليست مشتقة أيّ شيء أبدًا.
\end{example}

\begin{proposition}\label{prop:b1:continuity:algebra}
مجاميع الدوال \omterm{def:b1:continuity:continuous}{المتصلة} وجداءاتها وقسمتها (حيث تكون معرَّفة) وتراكيبها
\omterm{def:b1:continuity:continuous}{متصلة}. وكثيرات الحدود والكسور الناطقة
(خارج أقطابها) و $\abs{\,\cdot\,}$ و $\exp$ و $\ln$ والدوال
المثلثية والزائدية ومعكوساتها
(\cref{ch:b1:functions}) \omterm{def:b1:continuity:continuous}{متصلة} على مجموعات تعريفها.
\end{proposition}

\begin{proof}
العمليات: \cref{cor:b1:continuity:operations}. وأمّا الثوابت و
\omterm{def:b1:logic:map}{التطبيق} المطابق \omterm{def:b1:continuity:continuous}{فمتصلة} مباشرةً من التعريف
(إذ يفي $\delta = \varepsilon$ \omterm{def:b1:logic:map}{بالتطبيق} المطابق، ويفي أيّ $\delta$
بالثوابت)؛ ولأن جداءات الدوال \omterm{def:b1:continuity:continuous}{المتصلة} \omterm{def:b1:continuity:continuous}{متصلة}،
يتبع كل وحيد حدّ $a_k x^k$ بالاستقراء على $k$، وتُتمّ المجاميع
كثيرات الحدود؛ \omterm{def:b1:fractions:field}{والكسر الناطق} قسمة كثيرَي
حدود، وهو متصل حيثما لا ينعدم المقام.
وأمّا $\abs{\,\cdot\,}$: فبمتراجحة المثلث المعكوسة،
$\bigl|\abs{f(x)} - \abs{f(x_0)}\bigr| \leq \abs{f(x) -
f(x_0)}$، ومنه يفي $\delta$ نفسه بالغرض. وأمّا الدوال
الكلاسيكية فنمنحها \omterm{def:b1:continuity:continuous}{الاتصال} هنا؛ والاشتقاقية (المبرهن عليها في
\cref{ch:b1:derivative}) أقوى.
\end{proof}

\begin{example}[القيمة العظمى والصغرى لدالتين متصلتين]\label{ex:b1:continuity:maxmin}
إذا كانت $f, g$ متصلتين، فكذلك $\max(f, g)$ و $\min(f, g)$:
دون أيّ حاجة إلى تحليل حالات، بفضل المتطابقتين
\[
\max(f, g) = \frac{f + g + \abs{f - g}}{2},
\qquad
\min(f, g) = \frac{f + g - \abs{f - g}}{2},
\]
\omterm{def:b1:continuity:continuous}{واتصال} المجاميع \omterm{def:b1:continuity:continuous}{واتصال} $\abs{\,\cdot\,}$
(\cref{prop:b1:continuity:algebra}). وعلى الخصوص تكون $f^+ =
\max(f, 0)$ و $f^- = \max(-f, 0)$ متصلتين مع $f =
f^+ - f^-$: ومنه ففصل الإشارة المستعمل من أجل المتسلسلات
(\cref{ch:b1:series})، وعلى نطاق كامل في نظرية
التكامل في مجلّد السنة الثالثة، لا يكلّف شيئًا في الانتظام.
\end{example}

\begin{theorem}[مبرهنة القيم الوسطى]\label{thm:b1:continuity:ivt}
لتكن $f$ \omterm{def:b1:continuity:continuous}{متصلة} على $\intcc{a}{b}$ مع $f(a) \leq 0 \leq
f(b)$. عندئذ $f(c) = 0$ من أجل $c \in \intcc{a}{b}$ ما.
ونتيجةً لذلك، تأخذ دالة \omterm{def:b1:continuity:continuous}{متصلة} على \omterm{prop:b1:reals:intervals}{فترة} كلَّ قيمة
بين أيّ قيمتين لها: فتكون $f(I)$
\omterm{prop:b1:reals:intervals}{فترة}.\index{مبرهنة القيم الوسطى}
\end{theorem}

\begin{proof}
بالثنائية. ضع $a_0 = a$ و $b_0 = b$. ومن أجل $\intcc{a_n}{b_n}$ مع
$f(a_n) \leq 0 \leq f(b_n)$، ليكن $m$ المنتصف: فإذا كان $f(m) \leq
0$ فاحتفظ بالمقدار $\intcc{m}{b_n}$، وإلا فاحتفظ بالمقدار $\intcc{a_n}{m}$؛ فتدوم
شروط الإشارة. والمتتاليتان $(a_n), (b_n)$ متجاورتان ($b_n -
a_n = \frac{b-a}{2^n}$)، ونهايتهما المشتركة $c$
(\cref{thm:b1:seq:adjacent}). وبالاتصال و
\cref{thm:b1:seq:order}: $f(c) = \lim f(a_n) \leq 0$ و $f(c) =
\lim f(b_n) \geq 0$، ومنه $f(c) = 0$.

وأمّا النتيجة: فمن أجل قيم $f(u) < v < f(w)$، طبّق ما سبق
على $x \mapsto f(x) - v$ على القطعة ذات الطرفين $u$ و $w$؛
ومنه تكون $f(I)$ محدَّبة، أي \omterm{prop:b1:reals:intervals}{فترة}
(\cref{prop:b1:reals:intervals}).
\end{proof}

\begin{omfigure}
\begin{tikzpicture}
\begin{axis}[omaxis, width=8.4cm, height=5.6cm,
    xmin=0, xmax=4.4, ymin=-1.1, ymax=1.4,
    xtick={0.6,3.8}, xticklabels={$a$,$b$}, ytick=\empty]
  \addplot[omDef, very thick, domain=0.6:3.8, samples=160]
    {0.5*(x-2) + 0.4*sin(deg(2.5*x))};
  \draw[dashed, gray] (axis cs:0.6,0) -- (axis cs:0.6,-0.301);
  \draw[dashed, gray] (axis cs:3.8,0) -- (axis cs:3.8,0.870);
  \node[below, font=\small, omDef] at (axis cs:0.6,-0.32) {$f(a)$};
  \node[above, font=\small, omDef] at (axis cs:3.8,0.89) {$f(b)$};
\end{axis}
\end{tikzpicture}

{\small على دالة \omterm{def:b1:continuity:continuous}{متصلة} تحقق $f(a) < 0 < f(b)$ أن تقطع
المحور: فبرهان الثنائية في \cref{thm:b1:continuity:ivt} يحبس نقطة
قطع بين متتاليتين متجاورتين.}
\end{omfigure}

\begin{example}[معادلة واحدة، والبروتوكول كاملًا]\label{ex:b1:continuity:protocol}
حُلَّ $\eu^x = 3 - x$ على $\R$: الوجود والوحدانية
والموضع. ضع $g(x) = \eu^x + x - 3$، وهي \omterm{def:b1:continuity:continuous}{متصلة}. الموضع و
الوجود: $g(0) = -2 < 0$ و $g(1) = \eu - 2 > 0$، ومنه
تزرع مبرهنة القيم الوسطى حلًا في
$\intoo{0}{1}$. والوحدانية: $g$ مجموع الدالة المتزايدة قطعًا
$\eu^x$ والدالة $x - 3$، ومنه فهي متزايدة قطعًا على
$\R$؛ والدالة الرتيبة قطعًا تأخذ كل قيمة مرة واحدة
على الأكثر، ومنه فالحل وحيد على $\R$ كله (لا في
\omterm{prop:b1:reals:intervals}{الفترة} المسبورة فقط). والبروتوكول --- رتّب لتحصل على $g = 0$، ثم
تغيّر الإشارة من أجل الوجود، ثم الرتابة من أجل الوحدانية --- يفصل في
أكثر أسئلة «كم عدد الحلول» في ثلاثة أسطر، و
تمدّه جداول التغيرات في \cref{ch:b1:derivative}
إلى $g$ غير الرتيبة بقطع $\R$ إلى فروع
رتيبة.
\end{example}

\begin{example}[الثنائية خوارزمية]\label{ex:b1:continuity:dichotomy}
برهان \cref{thm:b1:continuity:ivt} يحسب. خذ $f(x) =
x^3 + x - 1$: $f(0) = -1 < 0 < 1 = f(1)$، ومنه يقع جذر في
$\intoo{0}{1}$. وبالتنصيف:
\[
f(0.5) = -0.375 < 0, \qquad
f(0.75) = 0.171875 > 0, \qquad
f(0.625) = -0.130859375 < 0 ,
\]
ومنه يُحبس الجذر تباعًا في $\intoo{0.5}{1}$، ثم
$\intoo{0.5}{0.75}$، ثم $\intoo{0.625}{0.75}$ (والقيمة الحقيقية:
$c \approx 0.6823$). وبعد $n$ خطوة يكون الخطأ على الأكثر
$\frac{b - a}{2^n}$: فعشر خطوات تعطي ثلاثة أرقام عشرية، وعشرون تعطي
ستة. والفكرة النافذة: مبرهنة القيم الوسطى ليست
\omterm{def:b1:logic:statement}{عبارة} وجود فحسب --- بل برهان الثنائية لها خوارزمية
مضمونة، وإن كانت بطيئة، لإيجاد الجذور، وينبغي أن تُقاس بها
طريقةُ نيوتن السريعة الموضعية في \cref{ch:b1:derivative}.
\end{example}

\begin{theorem}[مبرهنة القيم الحدّية]\label{thm:b1:continuity:evt}
كل دالة \omterm{def:b1:continuity:continuous}{متصلة} على قطعة $\intcc{a}{b}$ محدودة و
تبلغ حاصريها: إذ يوجد $c, d \in \intcc{a}{b}$ يحققان
\[
f(c) = \inf_{\intcc{a}{b}} f,
\qquad
f(d) = \sup_{\intcc{a}{b}} f .
\]
ومع \cref{thm:b1:continuity:ivt}: \emph{تكون الصورة \omterm{def:b1:continuity:continuous}{المتصلة}
لقطعة قطعةً} $\intcc{f(c)}{f(d)}$.
\index{مبرهنة القيم الحدّية}
\end{theorem}

\begin{proof}
\emph{محدودة من أعلى:} وإلا فاختر $u_n$ يحقق $f(u_n) \geq n$. وبتراص
القطعة (\cref{thm:b1:topology:compact})، تتقارب \omterm{def:b1:seq:subsequence}{متتالية
جزئية} $u_{\varphi(n)} \to x \in \intcc{a}{b}$؛ ويعطي \omterm{def:b1:continuity:continuous}{الاتصال}
$f(u_{\varphi(n)}) \to f(x)$، لكن $f(u_{\varphi(n)}) \geq \varphi(n)
\to +\infty$: وهذا تناقض.

\emph{بلوغ \omterm{def:b1:reals:bounds}{الحد الأعلى}:} ليكن $M = \sup f$ ولنختر $v_n$ يحقق $f(v_n) >
M - \frac{1}{n+1}$ (\cref{prop:b1:reals:epsilon}). واستخرج
$v_{\varphi(n)} \to d \in \intcc{a}{b}$: عندئذ $f(d) = \lim
f(v_{\varphi(n)}) = M$ بالحصر. ويُعالج \omterm{def:b1:reals:bounds}{الحد الأدنى} بالمقدار
$-f$.
\end{proof}

\begin{example}[أصغر عنصر موجب على قطعة]\label{ex:b1:continuity:posmin}
لتكن $f$ \omterm{def:b1:continuity:continuous}{متصلة} على $\intcc{0}{1}$ مع $f(x) > 0$ من أجل
كل $x$. عندئذ $\inf f > 0$: فحسب مبرهنة القيم الحدّية يكون
\omterm{def:b1:reals:bounds}{الحد الأدنى} \emph{قيمةً} $f(c)$، و $f(c) > 0$ بحكم الفرض.
ومنه فالدالة \omterm{def:b1:continuity:continuous}{المتصلة} الموجبة على قطعة تبقى بعيدة
عن $0$ --- وهي حجة من سطرين تُستعمل عشرات المرات في
الفصول الآتية (ضبط المقامات، وتأطيرات الدوال
الدرجية، وحواصر الأخطاء). وأمّا على \omterm{prop:b1:reals:intervals}{فترة} غير متراصة فهذا يفشل
فشلًا مذهلًا: فالدالة $f(x) = x$ على $\intoc{0}{1}$ \omterm{def:b1:continuity:continuous}{متصلة} و
موجبة مع $\inf f = 0$، غير مبلوغ. والفكرة النافذة:
ترتقي «موجبة» إلى «موجبة انتظامًا» بالضبط عندما تكون
\omterm{def:b1:logic:sets}{مجموعة} التعريف متراصة؛ فكل فرض من فروض مبرهنة القيم الحدّية
حامل.
\end{example}

\begin{remark}[مزالق شائعة حول المبرهنات الثلاث]
(أ) \emph{الصور \omterm{def:b1:continuity:continuous}{المتصلة}}: لا تصمد إلا \emph{القطع}.
فليس على الصورة \omterm{def:b1:continuity:continuous}{المتصلة} \omterm{prop:b1:reals:intervals}{لفترة} \omterm{def:b1:topology:open}{مفتوحة} أن تكون \omterm{def:b1:topology:open}{مفتوحة}
(إذ يرسل $x \mapsto x^2$ الفترةَ $\intoo{-1}{1}$ على $\intco{0}{1}$)، وليس على
صورة \omterm{def:b1:topology:closed}{مجموعة مغلقة} أن تكون \omterm{def:b1:topology:closed}{مغلقة} (إذ يرسل $\arctan$
المجموعةَ \omterm{def:b1:topology:closed}{المغلقة} $\R$ على \omterm{def:b1:topology:open}{المفتوحة}
$\intoo{-\frac\pi2}{\frac\pi2}$)؛ لكن صورة قطعة
قطعةٌ (\cref{thm:b1:continuity:evt}). (ب) \emph{تحتاج
مبرهنة القيم الوسطى إلى \omterm{prop:b1:reals:intervals}{فترة}}: فالدالة $x
\mapsto \frac1x$، \omterm{def:b1:continuity:continuous}{المتصلة} على $\intco{-1}{0} \cup
\intoc{0}{1}$، تأخذ القيمتين $-1$ و $1$ ومع ذلك لا تنعدم
أبدًا --- \omterm{def:b1:logic:sets}{فمجموعة} تعريفها قطعتان منفصلتان، وتقع القيمة
$0$ في الفجوة؛ فسمِّ دائمًا \omterm{prop:b1:reals:intervals}{الفترة} التي تُطبَّق عليها
المبرهنة. (ج) \emph{\omterm{def:b1:continuity:uniform}{الاتصال المنتظم} خاصية
للزوج (الدالة، \omterm{def:b1:logic:sets}{المجموعة})}: فالدالة $x^2$ \omterm{def:b1:continuity:uniform}{متصلة انتظامًا} على كل
قطعة ومع ذلك ليست كذلك على $\R$
(\cref{ex:b1:continuity:notuniform}) --- \omterm{def:b1:logic:statement}{فعبارة} «\omterm{def:b1:continuity:uniform}{متصلة
انتظامًا}» بلا \omterm{def:b1:logic:sets}{مجموعة} تعريف بلا معنى. (د)
\emph{\omterm{def:b1:continuity:continuous}{اتصال} المعكوس ليس صوريًا}: فهو يصحّ على
الفترات عبر الرتابة
(\cref{thm:b1:continuity:bijection})، لكن قد يكون لتقابل متصل
بين اتحادات فترات معكوسٌ غير متصل ---
والملاحظة بعد تلك المبرهنة موجودة لأن الطلبة يستشهدون بها
دون فرض \omterm{prop:b1:reals:intervals}{الفترة}.
\end{remark}

\section{الدوال الرتيبة والدوال المعكوسة}

\begin{theorem}[مبرهنة التقابل الرتيب]\label{thm:b1:continuity:bijection}
لتكن $f$ \omterm{def:b1:continuity:continuous}{متصلة} ورتيبة قطعًا على \omterm{prop:b1:reals:intervals}{فترة} $I$. عندئذ
\begin{enumerate}
  \item تكون $f$ تقابلًا من $I$ على \omterm{prop:b1:reals:intervals}{الفترة} $J = f(I)$؛
  \item ويكون المعكوس $f^{-1} \colon J \to I$ رتيبًا قطعًا (في
        الاتجاه نفسه) و\emph{متصلًا}.
\end{enumerate}
\end{theorem}

\begin{proof}
لنقل إن $f$ متزايدة قطعًا. (1) التباين مباشر من
الرتابة القطعية؛ والشمول على $f(I)$ بديهي، و $f(I)$
\omterm{prop:b1:reals:intervals}{فترة} حسب \cref{thm:b1:continuity:ivt}.

(2) الدالة $f^{-1}$ متزايدة قطعًا: فإذا كان $y < y'$ في $J$ لكن $f^{-1}(y)
\geq f^{-1}(y')$، فإن \omterm{def:b1:logic:map}{تطبيق} $f$ المتزايدة يعطي $y \geq y'$،
وهذا محال. وأمّا \omterm{def:b1:continuity:continuous}{اتصال} $f^{-1}$ عند $y_0 = f(x_0) \in J$: فليكن
$\varepsilon > 0$. افترض أولًا أن $x_0$ داخلية \omterm{def:b1:logic:sets}{للمجموعة} $I$، و
صغّر $\varepsilon$ بحيث يكون $x_0 \pm \varepsilon \in I$: فتحقق
صورهما $f(x_0 - \varepsilon) < y_0 < f(x_0 +
\varepsilon)$. خذ
$\delta = \min\bigl(y_0 - f(x_0 - \varepsilon),\, f(x_0 +
\varepsilon) - y_0\bigr) > 0$: فمن أجل $\abs{y - y_0} \leq \delta$،
تحصر رتابةُ $f^{-1}$ المقدارَ $f^{-1}(y)$ بين $x_0 -
\varepsilon$ و $x_0 + \varepsilon$. وإذا كانت $x_0$، مثلًا، الطرف
الأيسر \omterm{def:b1:logic:sets}{للمجموعة} $I$، فلا يتوفر إلا $x_0 + \varepsilon$: عندئذ
$y_0 = \min J$ (بالرتابة)، ويحقق كل $y \in J$ مع $y - y_0
\leq f(x_0 + \varepsilon) - y_0$ الشرطَ $x_0 \leq f^{-1}(y)
\leq x_0 + \varepsilon$، والتقدير من جهة واحدة هو بالضبط
\omterm{def:b1:continuity:continuous}{الاتصال} عند طرف؛ والطرف الأيمن متناظر.
(ولاحظ: لم يُستنتج \omterm{def:b1:continuity:continuous}{اتصال} $f^{-1}$
من \omterm{def:b1:continuity:continuous}{اتصال} $f$ بالتناظر --- بل الرتابة على \omterm{prop:b1:reals:intervals}{فترة} هي التي تقوم بالعمل.)
\end{proof}

\begin{remark}
هذه المبرهنة هي ما استعمله \cref{def:b1:functions:arc} و
\cref{prop:b1:functions:invhyp} صامتين: فالدوال $\arcsin$ و $\arccos$ و
$\arctan$ و $\operatorname{arsinh}$ و\dots \omterm{def:b1:continuity:continuous}{متصلة}. وتكملةً
(\cref{exo:b1:continuity:10}): تكون الدالة \omterm{def:b1:continuity:continuous}{المتصلة}
\emph{المتباينة} على \omterm{prop:b1:reals:intervals}{فترة} رتيبةً قطعًا تلقائيًا، ومنه فلا يكلّف
فرض الرتابة شيئًا.
\end{remark}

\begin{example}[الجذور من كل الرتب]\label{ex:b1:continuity:nthroot}
من أجل $n \in \N^*$، تكون الدالة $f(x) = x^n$ \omterm{def:b1:continuity:continuous}{متصلة} و
متزايدة قطعًا على $\intco{0}{+\infty}$، مع $f(0) = 0$ و
$f(x) \to +\infty$: فصورتها هي $\intco{0}{+\infty}$
(\cref{thm:b1:continuity:ivt} من أجل بنية \omterm{prop:b1:reals:intervals}{الفترة}). وعندئذ
تعطي مبرهنة التقابل الرتيب، دفعة واحدة، معكوسًا
\emph{متصلًا} متزايدًا قطعًا
\[
x \mapsto x^{1/n} \colon \intco{0}{+\infty} \to
\intco{0}{+\infty} :
\]
أي وجود الجذور من الرتبة $n$ ووحدانيتها واتصالها، دون أيّ
حساب. وقارن بالمقدار \cref{exo:b1:reals:12}، الذي بنى
$\sqrt y$ باليد من \omterm{def:b1:reals:bounds}{الحد الأعلى}: فقد ضغط فصلٌ واحد من النظرية
تلك الصفحة من العمل في سطرين، والسطران نفساهما
شرّعا $\arcsin$ و $\arctan$ و
$\operatorname{arsinh}$ قبلها. والفكرة النافذة: المبرهنة الجيدة
عملٌ مخزون.
\end{example}

\section{الاتصال المنتظم}

\begin{definition}\label{def:b1:continuity:uniform}
تكون $f \colon I \to \R$ \emph{متصلة انتظامًا}\index{اتصال
منتظم} إذا كان
\[
\forall \varepsilon > 0,\ \exists \delta > 0,\ \forall x, y \in I,
\qquad \abs{x - y} \leq \delta \implies \abs{f(x) - f(y)} \leq
\varepsilon .
\]
والمقصود: أن $\delta$ لا يتعلق إلا بالمقدار $\varepsilon$، لا
بالموضع في $I$. \omterm{def:b1:continuity:continuous}{والاتصال} المنتظم يستلزم \omterm{def:b1:continuity:continuous}{الاتصال}؛ وكل دالة
\emph{ليبشيتزية} ($\abs{f(x) - f(y)} \leq k \abs{x - y}$)
\omterm{def:b1:continuity:continuous}{متصلة} انتظامًا ($\delta = \varepsilon/k$).
\end{definition}

\begin{example}\label{ex:b1:continuity:notuniform}
الدالة $x \mapsto x^2$ \omterm{def:b1:continuity:continuous}{متصلة} على $\R$ لكنها \emph{ليست}
\omterm{def:b1:continuity:uniform}{متصلة انتظامًا}: $\abs{(n + \frac 1n)^2 - n^2} = 2 + \frac{1}{n^2} \geq 2$
مع أن المتغيرين على بعد $\frac 1n \to 0$. وأمّا على أيّ \omterm{prop:b1:reals:intervals}{فترة}
محدودة فهي ليبشيتزية، ومنه \omterm{def:b1:continuity:uniform}{متصلة انتظامًا} ---
وهذا متسق مع مبرهنة هاينه أدناه.
\end{example}

\begin{example}[معاملات الانتظام، صراحةً]\label{ex:b1:continuity:moduli}
على قطعة، يضمن هاينه وجود $\delta$ منتظم؛ وكثيرًا ما يمكن
\emph{حسابه} كذلك. فمن أجل $f(x) = x^2$ على $\intcc{0}{10}$:
\[
\abs{x^2 - y^2} = \abs{x + y}\,\abs{x - y} \leq 20\,\abs{x - y},
\]
ومنه يفي $\delta = \frac{\varepsilon}{20}$ بالغرض انتظامًا (وهو معامل
ليبشيتزي، خطي في $\varepsilon$). ومن أجل $\sqrt x$ على
$\intcc{0}{1}$: $\abs{\sqrt x - \sqrt y} \leq \sqrt{\abs{x -
y}}$ (\cref{exo:b1:continuity:9})، ومنه يفي $\delta = \varepsilon^2$
بالغرض --- وهو منتظم لكنه \emph{ليس} خطيًا: فقرب $0$ يكون الجذر
التربيعي شديد الانحدار، ويظهر الثمن في أُسّ
$\varepsilon$، لا في فشل الانتظام. والفكرة
النافذة: \omterm{def:b1:continuity:uniform}{الاتصال المنتظم} طيف، لا نعم/لا ---
فالدالة $\delta(\varepsilon)$، المسمّاة المعامل، تقيس
كم يكلّف الانتظام، وليبشيتز ببساطة أفضل درجاتها.
\end{example}

\begin{theorem}[هاينه]\label{thm:b1:continuity:heine}
كل دالة \omterm{def:b1:continuity:continuous}{متصلة} على \emph{قطعة} $\intcc{a}{b}$ \omterm{def:b1:continuity:uniform}{متصلة
انتظامًا}.\index{مبرهنة هاينه}
\end{theorem}

\begin{proof}
بالخلف: افترض أن $\varepsilon_0 > 0$ ما يهزم كل
$\delta$. ومع $\delta = \frac{1}{n+1}$، اختر $x_n, y_n \in
\intcc{a}{b}$ يحققان $\abs{x_n - y_n} \leq \frac{1}{n+1}$ و
$\abs{f(x_n) - f(y_n)} > \varepsilon_0$. ويستخرج التراص
(\cref{thm:b1:topology:compact}) المتتاليةَ $x_{\varphi(n)} \to c \in
\intcc{a}{b}$؛ عندئذ $y_{\varphi(n)} \to c$ كذلك (بالحصر على $\abs{x -
y}$). ويعطي \omterm{def:b1:continuity:continuous}{الاتصال} عند $c$ أن $f(x_{\varphi(n)}) \to f(c)$ و
$f(y_{\varphi(n)}) \to f(c)$، ومنه $\abs{f(x_{\varphi(n)}) -
f(y_{\varphi(n)})} \to 0 < \varepsilon_0$: وهذا تناقض.
\end{proof}

\begin{example}[محدودة ومتصلة ومع ذلك ليست منتظمة]\label{ex:b1:continuity:sinx2}
الدالة $f(x) = \sin(x^2)$ \omterm{def:b1:continuity:continuous}{متصلة} ومحدودة على
$\R$، لكنها ليست \omterm{def:b1:continuity:uniform}{متصلة انتظامًا}. خذ
\[
x_n = \sqrt{2\pi n}, \qquad y_n = \sqrt{2\pi n + \tfrac\pi2}:
\qquad
y_n - x_n = \frac{\pi/2}{\sqrt{2\pi n + \frac\pi2} +
\sqrt{2\pi n}} \longrightarrow 0 ,
\]
ومع ذلك $f(y_n) - f(x_n) = \sin\bigl(2\pi n + \frac\pi2\bigr) -
\sin(2\pi n) = 1 - 0 = 1$ من أجل كل $n$: فلا $\delta$ واحد يستطيع
أن يفي بالمقدار $\varepsilon = \frac12$ في كل مكان. وهندسيًا،
\emph{تتسارع} تذبذبات $\sin(x^2)$: إذ يُتمّ المنحنى
موجة كاملة على نوافذ أقصر فأقصر، ومنه فالسلّم
الأفقي الذي يقتضيه $\varepsilon$ معطى يتقلص إلى الصفر
كلما كبر $x$. والفكرة النافذة: الحدّية لا تشتري
الانتظام (وهذا المثال شاهده)، وعدم الحدّية لا يمنعه
($\sqrt x$، \cref{exo:b1:continuity:9})؛ وما يفصل هو
\emph{معامل التذبذب}، وتقول مبرهنة هاينه إن مجموعات التعريف
المتراصة تؤدّبه تلقائيًا.
\end{example}

\begin{example}[تجربة الآلة الحاسبة، مفسَّرة]\label{ex:b1:continuity:cositeration}
اطبع أيّ عدد في آلة حاسبة واضغط $\cos$ مرارًا:
فتستقر الشاشة على $0.7390851\dots$ فلماذا؟ بعد ضغطة واحدة تقع
القيمة في $\intcc{-1}{1}$، وبعد ضغطتين في $\intcc{\cos 1}{1}
\subseteq \intcc{0.54}{1}$، وهي \omterm{prop:b1:reals:intervals}{فترة} مستقرة من أجل $\cos$. وعليها،
$\abs{\cos a - \cos b} \leq \sin(1)\,\abs{a - b}$ مع $\sin 1 =
0.841\dots < 1$ (وهو حاصر تحويل الجداء إلى مجموع في
\cref{pb:b1:seq:1}، أو متراجحة التزايدات المنتهية في
\cref{ch:b1:derivative}): ومنه فالتكرار تقلّصي، ومنه، حسب
خطوة ضبط الخطأ في \cref{met:b1:seq:recurrent}،
\[
\abs{u_n - c} \leq (0.842)^{\,n-2}\,\abs{u_2 - c}
\longrightarrow 0 ,
\]
حيث $c$ هي النقطة الصامدة الوحيدة $\cos c = c$
(\cref{exo:b1:continuity:6}). ونحو $40$ ضغطة تشتري ثلاثة
أرقام عشرية ($0.842^{40} \approx 10^{-3}$) --- وهي سرعة هندسية،
أخفّ من سرعة الثنائية $2^{-n}$ لكل خطوة، لكن كل ضغطة تكلّف
ضربة مفتاح واحدة بينما تكلّف كل خطوة ثنائية تقييمَ إشارة
كاملًا. والفكرة النافذة: صورة النقطة الصامدة في
\cref{ch:b1:seq} ومبرهنات الوجود في هذا الفصل
نصفان لقصة واحدة --- فمبرهنة القيم الوسطى تجد $c$، والتقلّص يبلغه.
\end{example}

\begin{remark}[أين تعمل هذه المبرهنات تاليًا]
تحرّك كلٌّ من أعمدة هذا الفصل الثلاثة فصلًا لاحقًا. فمبرهنة
القيم الوسطى تغذّي كل حجة على وجود الحلول
ومبرهنةَ التقابل الرتيب؛ ومبرهنة القيم
الحدّية تحوّل مسائل الأمثلية إلى مبرهنات (فمبرهنة رول و
مبرهنة التزايدات المنتهية في \cref{ch:b1:derivative} تبدآن هناك
بالضبط)؛ ومبرهنة هاينه هي سبب إمكان مكاملة الدوال \omterm{def:b1:continuity:continuous}{المتصلة}
على القطع في \cref{ch:b1:integration} --- فالمقدار $\delta$
المنتظم هو ما يجعل مجاميع ريمان تتقارب. وفي مجلّد السنة
الثانية يعود الثلاثي نفسه في الفضاءات المعيارية، ويقوم
التراص بالعمل الذي تقوم به القطع هنا.
\end{remark}

\begin{remark}[منظورات داخل هذا المجلّد]
\omterm{def:b1:continuity:continuous}{الاتصال} على وشك أن يُتجاوز لكنه لا يتقاعد أبدًا.
فالفصل \Cref{ch:b1:derivative} يقوّيه إلى الاشتقاقية و
يردّ الجميل (فالاشتقاقية تستلزم \omterm{def:b1:continuity:continuous}{الاتصال})؛
و \cref{ch:b1:integration} يقوم عليه مرتين، عبر هاينه من أجل
البناء وعبر المبرهنة الأساسية، التي يرقّي
غرضها المركزي $x \mapsto \int_a^x f$ دالةً $f$ \omterm{def:b1:continuity:continuous}{متصلة} فحسب
إلى دالة أصلية من الصنف $C^1$. وفي
\cref{ch:b1:multivar}، يحمل \omterm{def:b1:continuity:continuous}{الاتصال} في متغيرين فخًّا
يستحق العرض المسبق: فالدالة $\frac{xy}{x^2 + y^2}$ (ممدَّدة
بالقيمة $0$) \omterm{def:b1:continuity:continuous}{متصلة} في $x$ من أجل كل $y$ مثبَّت وفي $y$ من أجل
كل $x$ مثبَّت، ومع ذلك ليست \omterm{def:b1:continuity:continuous}{متصلة} عند المبدأ --- إذ إنها على
القطر $x = y$ تساوي $\frac12$ باستمرار. \omterm{def:b1:continuity:continuous}{فالاتصال} المنفصل
أضعف قطعًا من \omterm{def:b1:continuity:continuous}{الاتصال}: فالتمييز التتالي
ينجو من الانتقال إلى $\R^2$، لكن يجب السماح للمتتاليات
بالاقتراب من \emph{كل} اتجاه، لا على المحاور فقط.
\end{remark}

\section{تمارين}

\begin{exercise}[$\star$]\label{exo:b1:continuity:1}
باستعمال التمييز التتالي، برهن على أن $x \mapsto
\sin\frac 1x$ لا نهاية له عند $0^+$ \emph{(أظهر متتاليتين)}.
وهل للمقدار $x \mapsto x \sin\frac 1x$ نهاية؟
\end{exercise}

\begin{exercise}[$\star$]\label{exo:b1:continuity:2}
ادرس \omterm{def:b1:continuity:continuous}{الاتصال} على $\R$ للمقدار $f(x) = \lfloor x \rfloor$ وللمقدار
$g(x) = x - \lfloor x \rfloor$، وللمقدار $h(x) = \lfloor x \rfloor +
(x - \lfloor x\rfloor)^2$.
\end{exercise}

\begin{exercise}[$\star$]\label{exo:b1:continuity:3}
برهن على أن للمعادلة $x^5 - 3x + 1 = 0$ ثلاثة حلول حقيقية على
الأقل \emph{(قوّم عند نقاط محسنة الاختيار وطبّق
\cref{thm:b1:continuity:ivt} على ثلاث قطع منفصلة)}.
\end{exercise}

\begin{exercise}[$\star$]\label{exo:b1:continuity:4}
برهن على أن لكل \omterm{def:b1:poly:def}{كثير حدود} من درجة فردية جذرًا حقيقيًا.
\end{exercise}

\begin{exercise}[$\star\star$]\label{exo:b1:continuity:5}
(نقطة صامدة) لتكن $f \colon \intcc{0}{1} \to \intcc{0}{1}$
\omterm{def:b1:continuity:continuous}{متصلة}. برهن على أن للمقدار $f$ نقطةً صامدة: أي $f(c) = c$ من أجل $c$
ما. وبيّن أنه لا يمكن إسقاط \omterm{def:b1:continuity:continuous}{الاتصال} ولا القطعة.
\end{exercise}

\begin{exercise}[$\star\star$]\label{exo:b1:continuity:6}
برهن على أن للمعادلة $\cos x = x$ حلًا حقيقيًا واحدًا بالضبط،
وعلى أنه يقع في $\intoo{0}{1}$.
\end{exercise}

\begin{exercise}[$\star\star$]\label{exo:b1:continuity:7}
لتكن $f \colon \R \to \R$ \omterm{def:b1:continuity:continuous}{متصلة} مع $f(x) \to +\infty$ عندما $x
\to \pm\infty$. برهن على أن $f$ تبلغ أصغر قيمة شاملة على $\R$.
\emph{(ردّ الأمر إلى قطعة تحتوي \omterm{def:b1:logic:sets}{مجموعة} تحت المستوى.)}
\end{exercise}

\begin{exercise}[$\star\star$]\label{exo:b1:continuity:8}
لتكن $f \colon \R \to \R$ \omterm{def:b1:continuity:continuous}{متصلة} ودورية (بالدور $T >
0$). برهن على أن $f$ محدودة وتبلغ حاصريها، وعلى أنه
يوجد $c$ يحقق $f(c + \frac T2) = f(c)$. \emph{(ومن أجل النقطة
الثانية، ادرس $g(x) = f(x + \frac T2) - f(x)$ على دور واحد.)}
\end{exercise}

\begin{exercise}[$\star\star$]\label{exo:b1:continuity:9}
برهن على أن $x \mapsto \sqrt x$ \omterm{def:b1:continuity:uniform}{متصلة انتظامًا} على
$\intco{0}{+\infty}$، وإن لم تكن ليبشيتزية بجوار $0$.
\emph{(برهن على $\abs{\sqrt x - \sqrt y} \leq \sqrt{\abs{x -
y}}$ واستعمله.)}
\end{exercise}

\begin{exercise}[$\star\star\star$]\label{exo:b1:continuity:10}
لتكن $f$ \omterm{def:b1:continuity:continuous}{متصلة} ومتباينة على \omterm{prop:b1:reals:intervals}{فترة} $I$. برهن على أن
$f$ رتيبة قطعًا. \emph{إرشاد: إذا لم تكن كذلك، فيوجد $a < b < c$
مع، مثلًا، $f(b) > f(a)$ و $f(b) > f(c)$؛ فطبّق مبرهنة القيم
الوسطى على قيمة بين $\max(f(a), f(c))$ و $f(b)$ على
جانبَي $b$ معًا.}
\end{exercise}

\begin{exercise}[$\star\star\star$]\label{exo:b1:continuity:11}
(معادلة كوشي الدالية، الحالة \omterm{def:b1:continuity:continuous}{المتصلة}) لتكن $f \colon \R \to
\R$ \omterm{def:b1:continuity:continuous}{متصلة} مع $f(x + y) = f(x) + f(y)$ من أجل كل $x, y$.
برهن على أن $f(x) = f(1)\,x$ من أجل كل $x$: أولًا على $\N$ و $\Z$ و $\Q$
(بالجمعية وحدها)، ثم على $\R$ \omterm{def:b1:continuity:continuous}{بالاتصال} \omterm{def:b1:topology:dense}{والكثافة}
(\cref{thm:b1:reals:density}).
\end{exercise}

\begin{exercise}[$\star\star\star$]\label{exo:b1:continuity:12}
لتكن $f \colon \intco{0}{+\infty} \to \R$ \omterm{def:b1:continuity:continuous}{متصلة} مع
$f(x) \to \ell \in \R$ عندما $x \to +\infty$. برهن على أن $f$
\omterm{def:b1:continuity:continuous}{متصلة} \emph{انتظامًا} على $\intco{0}{+\infty}$. \emph{(اقطع عند
$A$ كبير: هاينه على $\intcc{0}{A+1}$، والنهاية بعد $A$؛ واجعل
النظامين متداخلين.)}
\end{exercise}

\section{مسألة: معادلة كوشي الدالية وأخواتها}

\begin{problem}[{مسألة نهاية الأسبوع --- $f(x + y) = f(x) + f(y)$:
الانتظام يفرض الخطية، وصورة الوحوش}]
\label{pb:b1:continuity:1}
أيّ الدوال تحقق $f(x + y) = f(x) + f(y)$ من أجل كل $x,
y$ حقيقيين؟ طرح كوشي السؤال سنة 1821؛ والجواب نموذج فكري.
ويبيّن \Cref{exo:b1:continuity:11} أن مثل هذه الدالة \emph{الجمعية}
خطية على $\Q$ وأن \omterm{def:b1:continuity:continuous}{الاتصال} الكامل يفرض $f(x)
= cx$. وتشحذ هذه المسألة الفرض شحذًا هائلًا ---
\omterm{def:b1:continuity:continuous}{فالاتصال} عند نقطة \emph{واحدة}، أو الرتابة، أو مجرد
الحدّية على \omterm{prop:b1:reals:intervals}{فترة} صغيرة واحدة، كلٌّ منها يكفي --- ثم ترسم
صورة حلّ لا خطي مفترض (فمنحناه
يملأ المستوي)، وتحلّ المعادلات الأخوات التي تميّز
$\eu^{cx}$ و $c\ln x$ و $x^c$ و $cx^2$، وتُختم بمعادلة ينسن
والمبرهنة \emph{محدَّبة عند المنتصف $+$ \omterm{def:b1:continuity:continuous}{متصلة}
$\implies$ محدَّبة}.\index{معادلة كوشي
الدالية}\index{معادلة دالية} وفي كل ما يأتي، تعني \emph{الجمعية}:
$f(x + y) = f(x) + f(y)$ من أجل كل $x, y \in \R$.

\medskip
\noindent\textbf{الجزء 1 --- الخطية على $\Q$، ونقطة \omterm{def:b1:continuity:continuous}{اتصال}
واحدة.}
\begin{enumerate}
  \item لتكن $f$ جمعية. من \cref{exo:b1:continuity:11}،
        يكون $f(r) = f(1)\,r$ من أجل $r$ ناطق. برهن على \omterm{def:b1:logic:statement}{العبارة}
        الأدقّ المستعملة أدناه: من أجل \emph{كل} $x \in \R$ و $r
        \in \Q$، $f(rx) = r\,f(x)$ (أي إن $f$
        خطية على $\Q$).
  \item افترض أن الدالة الجمعية $f$ \omterm{def:b1:continuity:continuous}{متصلة} عند نقطة واحدة
        $x_0$. بيّن أن $f$ \omterm{def:b1:continuity:continuous}{متصلة} في كل مكان
        \emph{(احسب $f(x + h) - f(x)$ بدلالة $f(x_0 + h)
        - f(x_0)$)}، ومنه $f(x) = f(1)\,x$.
  \item بيّن أن الدالة الجمعية تتحدد بقصرها على أيّ
        \omterm{def:b1:structures:subgroup}{زمرة جزئية} \omterm{def:b1:topology:dense}{كثيفة}: فإذا تطابقت دالتان جمعيتان
        على $\Z + \sqrt2\,\Z$ وكانتا متصلتين، فهما
        متساويتان --- بينما، دون \omterm{def:b1:continuity:continuous}{اتصال}، يكون
        فرض $f(1) = 0$ و $f(\sqrt 2) = 1$ متسقًا
        مع الخطية على $\Q$ على \omterm{def:b1:structures:subgroup}{الزمرة الجزئية}. واحسب $f(m +
        n\sqrt2)$ من أجل هذا الفرض.
  \item لتكن $f$ جمعية ومحدودة من أعلى بالمقدار $M$ على \omterm{prop:b1:reals:intervals}{فترة} ما
        $\intcc{a}{b}$ مع $a < b$. بيّن أن $f$
        محدودة من أعلى على $\intcc{0}{\ell}$ حيث $\ell = b - a$
        \emph{(بالانسحاب بالمقدار $a$)}.
\end{enumerate}

\medskip
\noindent\textbf{الجزء 2 --- سلّم الانتظام.}
\begin{enumerate}[resume]
  \item متابعةً للسؤال 4: باستعمال $f(\ell - t) + f(t) =
        f(\ell)$، بيّن أن $f$ محدودة كذلك \emph{من أدنى} على
        $\intcc{0}{\ell}$: أي $\abs f \leq C$ هناك.
  \item بيّن $\abs{f(t)} \leq \frac{C}{n}$ من أجل $t \in
        \intcc{0}{\ell/n}$، واستنتج أن $f$ \omterm{def:b1:continuity:continuous}{متصلة} عند
        $0$ \emph{(والفردية تعالج الجهة اليسرى)}، ومنه
        في كل مكان (السؤال 2): فالدالة الجمعية المحدودة على
        \omterm{prop:b1:reals:intervals}{فترة} واحدة خطية.
  \item استنتج الحالة الرتيبة: الدالة الجمعية
        غير المتناقصة على $\intcc{a}{b}$ ما ($a < b$) هي $f(x)
        = cx$ مع $c \geq 0$.
  \item ركّب \emph{سلّم الانتظام}: من أجل $f$
        جمعية، تتكافأ العبارات الآتية --- (أ) $f(x) = cx$؛
        (ب) $f$ \omterm{def:b1:continuity:continuous}{متصلة}؛ (ج) $f$ \omterm{def:b1:continuity:continuous}{متصلة} عند نقطة واحدة؛ (د)
        $f$ رتيبة على \omterm{prop:b1:reals:intervals}{فترة} ما غير منحلّة؛ (ه) $f$
        محدودة على \omterm{prop:b1:reals:intervals}{فترة} ما غير منحلّة. ورتّب
        الاستلزامات بحيث يكون كلٌّ إمّا بديهيًا وإمّا مبرهنًا
        عليه أصلًا.
  \item تحقق من أن الأسئلة 4--6 لم تستهلك إلا حاصرًا
        \emph{من أعلى}: فالدالة الجمعية المحدودة من أعلى على \omterm{prop:b1:reals:intervals}{فترة}
        واحدة غير منحلّة خطية أصلًا. واستنتج
        \omterm{def:b1:logic:statement}{العبارة} المرآتية من أجل حاصر من أدنى، وسجّل
        أقوى صيغة لدرجة السلّم (ه) محصَّلة هكذا.
\end{enumerate}

\medskip
\noindent\textbf{الجزء 3 --- صورة وحش.} افترض
الآن أن $f$ جمعية لكنها \emph{ليست} خطية.
\begin{enumerate}[resume]
  \item بيّن أنه يوجد عددان حقيقيان غير معدومين $u, v$ يحققان
        $\dfrac{f(u)}{u} \neq \dfrac{f(v)}{v}$، وأن
        المتجهتين $(u, f(u))$ و $(v, f(v))$ يولّدان المستوي (فمحددهما
        $u f(v) - v f(u)$ غير معدوم).
  \item بيّن أن منحنى $f$ يحتوي كل النقاط
        \[
        r\,(u, f(u)) + s\,(v, f(v)), \qquad r, s \in \Q ,
        \]
        واستنتج أن المنحنى \emph{كثيف في $\R^2$}: أي إنه من أجل
        كل نقطة $(x_0, y_0)$ من المستوي وكل
        $\varepsilon > 0$، تقع $(x, f(x))$ ما على بعد
        $\varepsilon$ منها \emph{(حُلَّ الجملة الحقيقية
        $2\times2$، ثم قارب المعاملات الحقيقية بأعداد
        ناطقة)}.
  \item استنتج من السؤال 11 الصورة الكاملة: الدالة الجمعية
        غير الخطية غير محدودة على كل \omterm{prop:b1:reals:intervals}{فترة} غير منحلّة،
        ومنقطعة عند كل نقطة، وغير رتيبة على أيّ
        \omterm{prop:b1:reals:intervals}{فترة}، وصورتها لأيّ \omterm{prop:b1:reals:intervals}{فترة} \omterm{def:b1:topology:dense}{كثيفة} في
        $\R$. ووفّق ذلك مع السؤال 8.
  \item والوحوش موجودة --- على \omterm{def:b1:structures:subgroup}{زمرة جزئية} \omterm{def:b1:topology:dense}{كثيفة}، بناءً:
        على $G = \Z + \sqrt2\,\Z$ عرّف $f(m + n\sqrt2) = n$.
        بيّن أن $f$ معرَّفة جيدًا وجمعية على $G$، وأن
        $f$ غير محدودة على $G \cap \intoo{0}{\varepsilon}$ من أجل
        كل $\varepsilon > 0$ \emph{(فمن أجل $N$ مثبَّت، لا يوجد إلا
        عدد منته من $g = m + n\sqrt2 \in \intoo{0}{1}$ يحقق
        $\abs n \leq N$؛ ومع ذلك $G \cap \intoo{0}{\varepsilon}$
        غير منتهية)}. وفسّر في فقرة واحدة لماذا يقتضي تمديد
        مثل هذه الدالة $f$ إلى $\R$ كله أساسًا للمقدار $\R$ بوصفه
        فضاءً متجهيًا على $\Q$ (أي \emph{أساس هامل})، وهو
        وجودٌ مسألته من مسائل بديهية الاختيار وخارج هذا
        المجلّد.
\end{enumerate}

\medskip
\noindent\textbf{الجزء 4 --- المعادلات الأخوات.} وكل
الدوال هنا \omterm{def:b1:continuity:continuous}{متصلة}.
\begin{enumerate}[resume]
  \item لتكن $f \colon \R \to \R$ \omterm{def:b1:continuity:continuous}{متصلة} وغير مساوية $0$
        تمامًا، مع $f(x + y) = f(x)f(y)$. بيّن $f(x) =
        f\bigl(\frac x2\bigr)^2 \geq 0$، ثم $f > 0$ في كل مكان،
        ثم $f(x) = \eu^{cx}$ من أجل $c$ ما: فالأسّيات
        هي بالضبط التشاكلات \omterm{def:b1:continuity:continuous}{المتصلة} من $(\R, +)$ إلى
        $(\R^*, \times)$.
  \item لتكن $f \colon \intoo{0}{+\infty} \to \R$ \omterm{def:b1:continuity:continuous}{متصلة}
        مع $f(xy) = f(x) + f(y)$. بيّن $f(x) = c\ln x$
        \emph{(بالنقل بالمقدار $\exp$)}.
  \item لتكن $f \colon \intoo{0}{+\infty} \to \intoo{0}{+\infty}$
        \omterm{def:b1:continuity:continuous}{متصلة} مع $f(xy) = f(x)f(y)$. بيّن $f(x) = x^c$.
  \item جد كل الدوال \omterm{def:b1:continuity:continuous}{المتصلة} $f \colon \R \to \R$ التي تحقق
        \[
        f(x + y) = f(x) + f(y) + f(x)f(y) .
        \]
        \emph{(ادرس $h = 1 + f$؛ وعالج الحالة المنحلّة
        على حدة.)}
  \item (متوازي الأضلاع) جد كل الدوال \omterm{def:b1:continuity:continuous}{المتصلة} $f \colon \R \to \R$
        التي تحقق $f(x + y) + f(x - y) = 2f(x) + 2f(y)$: بيّن أن $f$
        زوجية، و $f(0) = 0$، و $f(nx) = n^2 f(x)$ بالاستقراء، ثم
        $f(x) = f(1)\,x^2$. (وهذه المعادلة بصمة
        الأشكال التربيعية --- وهي قانون متوازي الأضلاع الذي يكشف،
        في مجلّد السنة الثانية، أيّ المعايير يأتي من جداء
        سلّمي.)
\end{enumerate}

\medskip
\noindent\textbf{الجزء 5 --- ينسن والتحدّب عند المنتصف.}
\begin{enumerate}[resume]
  \item (معادلة ينسن) لتكن $f \colon \R \to \R$
        \omterm{def:b1:continuity:continuous}{متصلة} مع $f\bigl(\frac{x+y}{2}\bigr) = \frac{f(x)
        + f(y)}{2}$. بيّن أن $g = f - f(0)$ تحقق
        $g\bigl(\frac x2\bigr) = \frac{g(x)}{2}$، واستنتج أن
        $g$ جمعية، واختم بأن $f(x) = cx + d$.
  \item افترض الآن \emph{المتراجحة} فقط: $f$ \omterm{def:b1:continuity:continuous}{متصلة}
        مع
        \[
        f\Bigl(\frac{x + y}{2}\Bigr) \leq \frac{f(x) +
        f(y)}{2} \qquad (x, y \in \R).
        \]
        برهن بالاستقراء على $n$ على أنه من أجل كل الأوزان الثنائية
        $\lambda = \frac{k}{2^n} \in \intcc{0}{1}$:
        \[
        f\bigl(\lambda x + (1 - \lambda)y\bigr) \leq \lambda
        f(x) + (1 - \lambda) f(y) .
        \]
  \item مدّد \omterm{def:b1:continuity:continuous}{بالاتصال} وبكثافة الأعداد الثنائية
        (\cref{exo:b1:reals:8}) إلى كل $\lambda \in
        \intcc{0}{1}$: فالدالة \omterm{def:b1:continuity:continuous}{المتصلة} المحدَّبة عند المنتصف
        تحقق متراجحة التحدّب الكاملة (وهو المفهوم
        المدروس منهجيًا في \cref{ch:b1:derivative}).
  \item بيّن أنه لا يمكن إسقاط \omterm{def:b1:continuity:continuous}{الاتصال}: فالدالة الجمعية
        غير الخطية $f$ تحقق \emph{مساواة المنتصف} في السؤال
        19 ومع ذلك لا تحقق أيّ متراجحة تحدّب على أيّ \omterm{prop:b1:reals:intervals}{فترة} (السؤال
        12). والعبرة: التحدّب عند المنتصف خاصية ذات مراحل
        قابلة للعدّ (الأعداد الثنائية)، والتحدّب خاصية \omterm{def:b1:continuity:continuous}{متصلة}؛ و
        \omterm{def:b1:continuity:continuous}{الاتصال} هو الجسر --- تمامًا كما في الجزأين 1--2.
\end{enumerate}

\medskip
\noindent\textbf{الجزء 6 --- تنويعات أخيرة وتوليفة.}
\begin{enumerate}[resume]
  \item جد كل الدوال \omterm{def:b1:continuity:continuous}{المتصلة} $f \colon \R \to \R$ التي تحقق $f(x + y)
        = f(x) + f(y) + xy$ \emph{(اطرح الحل
        الخاص $\frac{x^2}{2}$)}.
  \item برهن: إذا كانت $f \colon \R \to \R$ \omterm{def:b1:continuity:continuous}{متصلة} و
        جمعية على \emph{\omterm{def:b1:structures:subgroup}{زمرة جزئية} \omterm{def:b1:topology:dense}{كثيفة}} $G$ فحسب (أي
        $f(g + g') = f(g) + f(g')$ من أجل $g, g' \in G$)، فإن $f$
        جمعية على $\R$. وأعمّ من ذلك، تتساوى دالتان متصلتان
        تتطابقان على \omterm{def:b1:logic:sets}{مجموعة} جزئية \omterm{def:b1:topology:dense}{كثيفة} من $\R$.
  \item توليفة، جملة واحدة لكلٍّ: (أ) اذكر سلّم
        الانتظام من السؤال 8 عن ظهر قلب؛ (ب) فسّر لماذا
        تكون «منحنى كثيف في المستوي» هي الصورة الذهنية
        الصحيحة لفشل الانتظام؛ (ج) اذكر الدوال الكلاسيكية
        الخمس المميَّزة في الجزأين 4--5 والطريقة
        الوحيدة التي أمسكت بها كلها؛ (د) وسمِّ الموضعين
        اللذين حملت فيهما \omterm{def:b1:topology:dense}{كثافة} $\Q$ (أو الأعداد الثنائية) في $\R$
        الحجةَ، والموضع الذي لم تستطع فيه (السؤال 13).
\end{enumerate}
\end{problem}
